\documentclass[12pt]{article}
\usepackage[utf8]{inputenc}
\usepackage[T1]{fontenc}
\usepackage{hyperref}
\title{Uruguayan Sign Language Language Resources for K-12 Curriculum Development}
\author{Compiled by Claude Sonnet 5 (claude-sonnet-5) for Andr\'as Kornai}
\date{}
\begin{document}
\maketitle

\section{Introduction}
Uruguayan Sign Language (Lengua de Se\~nas Uruguaya, LSU; ISO 639-3 ugy) was recognised by Law 17.378 of 25 July 2001 as the natural language of deaf people and their communities. This report rests on one Spanish-language search that returned mostly academic literature.

\section{Primary Educational Resources}
No LSU textbook or curriculum was found online. A study by Leonardo Peluso Crespi (Universidad de la Rep\'ublica) describes three models in public education (bilingual, bilingual-inclusive, inclusive) and proposes a plurilingual and intercultural model. It is among the sources at:

\url{https://www.redalyc.org/journal/3131/313169978059/313169978059.pdf}

\url{https://www.redalyc.org/journal/3131/313158902088/313158902088.pdf}

Which article is which was not verified.

\section{Linguistic and Reference Materials}
The \emph{Diccionario Biling\"ue de Lengua de Se\~nas Uruguaya/Espa\~nol} (ASUR/CINDE, Montevideo, 2007) appears in a university bibliography. No direct link to it was found. Related FHCE (Universidad de la Rep\'ublica) documents, including materials for the LSU--Spanish interpreting programme:

\url{https://www.fhce.edu.uy/images/Administracion_de_la_ensenanza/IMPAR\%202022/TUILSU/Estudios_Sordos.pdf}

\url{https://fhce.edu.uy/wp-content/uploads/2025/08/AES-web.pdf}

\section{Religious and Community Resources}
ASUR (Asociaci\'on de Sordos del Uruguay) is named as the dictionary's publisher. No URL was found.

\section{Digital and Audio Resources}
No video or corpus was found in the searches.

\section{Curriculum Development Context}
The legal basis is Law 17.378. Bilingual deaf classes divide the curriculum between LSU and Spanish. A 2010 figure cited in one work gave four deaf schools nationwide, which may be outdated. ASUR and FHCE are the first contacts for follow-up.
\end{document}